\section{Finite consequences of asymptotic estimates}
\label{sec:asymptotic-comparisons}

To decide whether an asymptotic result already gives
Theorem~\ref{thm:main}, we examine its finite-parameter estimates.
Write
\[
 \delta(s)=\dim_{\Q}\operatorname{span}_{\Q}
             \{1,\beta(2),\beta(4),\ldots,\beta(s)\}
 \qquad(s\text{ even}).
\]
A bound forcing $\delta(364)\geq4$ would imply three irrational values.
The comparisons below concern the specified estimates and parameter
families in the cited proofs. They give either a dimension bound too
small to force $4$, even after integer rounding, or a size bound that
does not establish decay of the integer linear forms.

\subsection{Rivoal--Zudilin's uniform family}

For the parameters in~\cite[(26)]{RZ03}, put $q=366$ and let
$r=2,4,\ldots,182$. Equation~(29) of that paper gives
\[
 \delta(q-2)\geq
 \frac{q\log r+(q-2r-1)\log2}
 {(2+\log2)(q-1)+(2r+1)\log(2r+1)}.
\]
The largest right-hand side is approximately $1.18156$, at $r=28$.
The sharper estimate preceding it also remains below $2$.
Indeed, evaluate the maximum defining its Laplace factor at
$t_1=\cdots=t_q=r/(r+1)$ and use
$(1+t_1\cdots t_q)^{2r+1}\leq2^{2r+1}$.
This bounds that sharper dimension estimate \emph{above} by
\begin{equation}\label{eq:rz-comparison}
 U(q,r)=\frac{q\{\log2+(r+1)\log(r+1)-r\log r\}}
 {(2+\log2)(q-1)+(2r+1)\log(2r+1)}<2.
\end{equation}
For all $91$ choices of $r$, the final inequality follows from the exact
integer comparison
\[
 8^{4(q-1)}2^{q-2}(2r+1)^{2(2r+1)}r^{qr}
 >3^{4(q-1)}(r+1)^{q(r+1)}
\]
and $\ep>8/3$. Also $U(q,r)$ increases with $q$ for fixed $r\geq2$,
so the same conclusion holds at smaller endpoints in this family.
These estimates can certify $\delta(364)\geq2$, hence one irrational
value, but cannot force dimension~$4$.

\subsection{Nishimoto's explicit estimate}

In~\cite[Theorem~5.5]{N11}, take the periodic coefficients to be $\ch$,
with period $d$ a positive multiple of $4$. For positive integers $a,b$
with $a$ even, put $r=(d+2b)/d$ and
$\Delta=(r+1)\log(r+1)-(r-1)\log(r-1)$.
The explicit estimate is $\delta(a)\geq1+\mathcal A/\mathcal B$, where
\begin{align*}
 \mathcal A&=(a+d)\Delta-2a(d+\log2)
       -4b(\log d-\log2)-2dr\log(2r)-\tfrac13,\\
 \mathcal B&=2ad+2a\log2+4(b+d)\log(b+d)-4d\log d.
\end{align*}
Among the hypotheses are $r\geq2$ and $(a+d)/d\geq3r$.
For $a\leq364$ these imply $a\geq5d$ and $r+1<32$, so
\[
 \Delta<2\log(r+1)+2<10\log2+2<9.
\]
Using $\log2>1/2$ and discarding negative terms gives
$\mathcal A<(8-2d)a+9d$.
For $d=4$, therefore, $\mathcal A<36$ and
$\mathcal B>9a\geq180$, whence $1+\mathcal A/\mathcal B<6/5$.
For $d\geq8$, we instead obtain
$\mathcal A<d(49-10d)<0$.
Thus this explicit estimate cannot force dimension~$4$ at~$364$.
For comparison, the period-$4$ entries in~\cite[Table~5]{N11} certify
dimension at least $2$ at $2594$ and at least $3$ at $990206$.

\subsection{Fischler's size estimates}

The estimates in~\cite[Section~4.2 and Proposition~4.4]{F18}
and~\cite[Proposition~1, (1.1)]{F20} bound the integer linear forms by
$\alpha^{n+o(n)}$. Their applications require $\alpha<1$.
With the subscripts distinguishing the two papers, the explicit factors are
\begin{align*}
 \alpha_{18}&=\ep^a4^{a/N-r}(N+1)^{2r+2}r^{-a/N+4r+2},
       &1\leq r&<\frac{a}{2N},\\
 \alpha_{20}&=(4\ep)^{(a+1)/N}(2N)^{2r+2}r^{-(a+1)/N+4(r+1)},
       &1\leq r&<\frac{a}{3N}.
\end{align*}
Here $a,N,r$ are positive integers. For even beta values we may take
$a=365$, since the largest even index is then $364$.
In the first construction, the ordinary representation of $\ch$ uses
$4\mid N$, while its twisted representation permits $N\equiv2\pmod4$,
including $N=2$. We therefore check all positive even $N$.
In the second construction we check all positive multiples of $4$,
including the enlarged periods used in the irrational-value count.

There are $429$ admissible pairs $(N,r)$ in the first case and $111$ in
the second. Every pair gives $\alpha>1$. These are exact checks:
raise $\alpha$ to the power $N$, replace $\ep$ by $8/3$, and compare
the resulting positive integers after clearing denominators.
For fixed $N,r$, each $\log\alpha$ is affine in $a$ and positive at
$a=0$. Positivity at $a=365$ therefore also covers smaller endpoints.
Thus these size estimates do not establish the decay needed at~$364$.

\subsection{The square-root bound}

For our result in~\cite{T26}, we check the finite dimension estimate
before the asymptotic choice of parameters. The argument below covers
all admissible choices at this endpoint. In that paper's notation, put $d=a+h$ and
$H_d=\sum_{j=1}^d1/j$. The largest beta index is
$2\lfloor d/2\rfloor$, so an endpoint at most $364$ means $d\leq365$.
Let $X$ be the coefficient-growth exponent $\log\chi$ in that construction.
Equations~(2.2), (5.2) and (5.3) of~\cite{T26} give
\begin{align*}
 X&\geq\frac{\Omega^2\log r}{2(a-\omega)},\\
 \log\alpha&=X-\Omega\log r+\kappa(1+\log4+H_d),\\
 \log B&=X+(r+1)\log(3/2)+\kappa(\log8+H_d).
\end{align*}
Here $a,h$ are integers with $h\geq0$, $r\geq2$, and
$a>\Omega\geq\omega>\kappa>2r$.
The dimension estimate is $\delta(2\lfloor d/2\rfloor)\geq1+R$,
where $R=-\log\alpha/\log B$ when $\alpha<1$.
We show that $R<2$, so this estimate cannot force dimension~$4$.

If $R\geq2$, rearranging the last two formulas yields
\[
 \Omega\log r-3X\geq
 \kappa(1+\log4+2\log8+3H_d)+2(r+1)\log(3/2).
\]
The lower bound on $X$ makes the left side at most the maximum of a
quadratic in $\Omega$, namely $(a-\omega)\log r/6$.
Since $\kappa>2r$ and $\log r/r\leq1/\ep$, division by $\kappa$ gives
\[
 \frac{\Omega\log r-3X}{\kappa}
 <\frac{365}{12\ep}<\frac{365}{32}.
\]
On the other hand, $d\geq5$, $\log4>1$ and $\log8>2$ make the right
side, divided by $\kappa$, greater than
\[
 6+3H_5=\frac{257}{20}>\frac{365}{32},
\]
a contradiction.

The accompanying script \texttt{checks/check\_asymptotic\_comparisons.py}
reproduces the finite integer comparisons and the rational constants used
above. It uses only Python's standard library; its decimal summaries
are illustrative and are not used to decide any inequality.
